\documentclass[11pt]{article}
\usepackage[utf8]{inputenc}
\usepackage[T1]{fontenc}
\usepackage{lmodern}
\usepackage[margin=1in]{geometry}
\usepackage{csquotes}
\usepackage[style=authoryear,backend=biber,sorting=nyt,maxbibnames=99,giveninits=true,uniquename=init]{biblatex}
\usepackage[hidelinks]{hyperref}
\addbibresource{references.bib}
\begin{document}
\section*{Conclusions}

Across observational, tracer, geophysical, experimental, and modelling lines
of evidence, a coherent and strongly supported picture of agricultural
subsurface drainage now emerges. Its foundation is hydrological: tile networks
are the master control on the water balance of drained land. By intercepting
water that would otherwise recharge shallow groundwater or leave as surface
runoff and routing it rapidly to streams, drains commonly carry $40$--$65\%$ of
streamflow in humid settings and dominate field-scale water export
\parencite{Schilling2019}. This
repartitioning reshapes the flow regime in consistent ways---elevating low
flows and the baseflow fraction, exerting only modest and scale-dependent
effects on the largest peaks, shortening groundwater travel times, and
homogenising catchment response so that drained basins behave far more
uniformly and scale-invariantly than their undrained counterparts. Relative to
undrained or surface-drained land, improved subsurface drainage generally
reduces surface runoff, sediment, and phosphorus even as it increases nitrate
delivery. Because this hydrology governs everything downstream, the physical
behaviour of drainage is the necessary starting point for understanding its
water-quality and agronomic consequences.

That hydrological primacy is clearest for nitrate, the dominant nitrogen
species leaving drained cropland and the best-quantified consequence of
drainage. The subsurface pathway carries the great majority of nitrogen loss;
concentrations routinely approach or exceed the $10~\mathrm{mg~L^{-1}}$
drinking-water limit; and loads are governed first by drainage discharge, so
export is strongly seasonal, concentrating in the non-growing season and in a
few high-flow episodes \parencite{Kladivko2021,Speir2021}. Concentration--discharge behaviour is broadly
chemostatic, and the nitrate is largely organic-derived and delivered rapidly
through macropores, with denitrification providing real but only partial
attenuation. Phosphorus tells a parallel but distinct story: although it moves
in far smaller mass, the long-standing assumption of a filtered ``clean''
subsurface pathway has been overturned. A tiny, hydraulically active volume of
biopores and cracks delivers particulate phosphorus and sediment to drains
within hours, so tiles frequently dominate edge-of-field phosphorus loads
despite lower concentrations than surface runoff, with dissolved losses tightly
coupled to soil phosphorus status and legacy phosphorus sustaining loads at
watershed scale \parencite{Pease2018,Sharpley1995}.

Against these losses the evidence supports a coherent, tiered hierarchy of
mitigation with quantified effects. In-field, nitrogen-rate optimisation toward
the economic optimum, cover crops---especially cereal rye, clustering around
$40$--$60\%$ nitrate-load reduction \parencite{Malone2018,Malone2014b,Malone2023}---diversified rotations, and perennials act
in predictable directions, while tillage is a weak and inconsistent lever. At
the design scale, spacing and depth trade drainage volume against loss, and
controlled drainage and drainage water recycling reduce loads largely in
proportion to the flow they retain, delivering water-quality benefits together
with yield resilience. At the edge of field, constructed wetlands, denitrifying
woodchip bioreactors, reactive phosphorus media, and saturated buffers
intercept loads through well-characterised controls---residence time,
temperature, loading, and pH \parencite{Christianson2013,Jaynes2008}. Two cross-cutting lessons recur: these practices
redistribute nitrogen and can trade dissolved for particulate phosphorus or one
loss pathway for another, so single-pathway evaluation misleads; and their
benefits are scale-attenuated, diluted between the tile and the watershed
outlet. Meeting downstream targets therefore requires stacking in-field and
edge-of-field measures, and traditional surface-water practices that fail to
intercept subsurface flow are largely ineffective over drained land. The same
preferential-flow and hydrologic controls extend to pesticides, manure-borne
microbes and resistance genes, trace elements, and dissolved greenhouse gases,
where transient concentrations can be high even as mass losses stay small and
management often matters more than drainage itself.

Alongside these environmental costs, drainage is an established and valuable
production technology. It reliably raises crop yields---most strongly on poorly
drained soils and in wet years---leaches salts to reclaim saline and
waterlogged land, and enables safe reuse of saline and brackish water where salt
balance is maintained. Controlled variants add further yield and water-use
benefits, and formal appraisals across many regions confirm favourable
financial and social returns, with careful design and mechanisation
substantially reducing cost. Underpinning the whole synthesis is a mature,
validated toolkit: process models such as \textsc{DRAINMOD} reproduce drainage
flow and nitrate reliably after calibration, with skill improving from daily to
annual scale ; the physics of saturated--unsaturated and preferential flow is
well understood; analytical spacing solutions, data-driven surrogates, and
remote sensing complement one another; and geophysical and machine-learning
methods now map buried networks and drained extent at field to national scales.

The strength of these conclusions rests on their convergence. Independent
methods, disparate climates, and separate scales point to the same mechanisms
and comparable magnitudes, so what is known about tile drainage is known with
some confidence. The unifying message is that drainage is fundamentally a
hydrological intervention: it sets the pathways, timing, and volume of water
movement, and water-quality outcomes, agronomic gains, and the effectiveness of
every management practice follow from that control. Recognising drainage water
quality as a design criterion alongside production, and combining complementary
practices across the field-to-watershed continuum, is the consolidated
implication of the evidence assembled here.

\printbibliography
\end{document}
